% Generated draft: requires mathematical and source review.
\documentclass[11pt]{article}
\usepackage[margin=25mm]{geometry}
\usepackage{fontspec}
\setmainfont{Latin Modern Roman}
\usepackage{amsmath,amssymb,mathtools,unicode-math}
\setmathfont{Latin Modern Math}
\usepackage{enumitem,needspace,xurl,hyperref}
\setlength{\parindent}{0pt}
\setlength{\parskip}{5pt}
\setlength{\emergencystretch}{3em}
\newcommand{\N}{\mathbb{N}}
\newcommand{\Z}{\mathbb{Z}}
\newcommand{\Q}{\mathbb{Q}}
\newcommand{\R}{\mathbb{R}}
\newcommand{\C}{\mathbb{C}}
\newcommand{\F}{\mathbb{F}}
\newcommand{\A}{\mathbb{A}}
\newcommand{\PP}{\mathbb{P}}
\newcommand{\E}{\mathbb{E}}
\newcommand{\eps}{\varepsilon}
\newcommand{\dd}{\,\mathrm{d}}
\begin{document}
\section*{UnsolvedMath ID: 20003275}
\textbf{Source problem code:} catalogue entry
\subsection{Definitions and mathematical statement}
Let $H$ act continuously on a space $X$, with a specified subgroup $P\leq H$ and standard Borel/Radon measure conventions. Call the action $P$-stiff when every finite positive $P$-invariant Borel measure is $H$-invariant. The source asks for geometric conditions ensuring stiffness and suggests the $SL(2,\R)$ action on quadratic-differential moduli spaces. It does not define $H,P$ or the full category of geometric actions in this isolated record.

For the supplied motivating case, take $H=SL(2,\R)$ and $P=AN$, the connected upper-triangular subgroup with positive diagonal. The aid attributes to Eskin--Mirzakhani that every ergodic $P$-invariant probability on a unit-area Abelian-differential stratum is $SL(2,\R)$-invariant and affine in period coordinates. It reports the corresponding quadratic-stratum invariance via the orientation double cover, with area normalization and deck-involution data treated compatibly. Ergodic decomposition then gives $H$-invariance for arbitrary finite nonergodic measures in those stated strata. The exact theorem references and transfer details are not supplied; this is not an unqualified assertion about every unspecified geometric action or moduli compactification.

The aid's finite-root diagnostic can be stated for a connected split semisimple real Lie group $H$, a fixed simple-root system $\Delta$, and a standard parabolic $P$ whose Levi simple roots are $\Theta\subset\Delta$. If a finite measure $\mu$ is $P$-invariant and additionally invariant under each negative simple-root subgroup $U_{-\alpha}$ for $\alpha\in\Delta\setminus\Theta$, then it is $H$-invariant, since these groups together generate $H$. The split/connected hypothesis makes this formulation explicit; other real forms require the corresponding generation theorem. This is a sufficient measure-symmetry test, not a claim that $P$-invariance automatically supplies those extra symmetries.

For $SL(2,\R)$ the aid gives the sharper test: if $\mu$ is $AN$-invariant and $g_*\mu=\mu$ for one $g\notin N_H(AN)$, then $\mu$ is $H$-invariant. The group generated by $AN$ and such a symmetry is dense in $H$, and the stabilizer of a finite measure is closed for the continuous action. A central element or a different component of the normalizer does not suffice.

The factor observation has an essential compactness condition. If $P$ is amenable and $\pi:X\twoheadrightarrow Y$ is a continuous equivariant map of compact Hausdorff $H$-spaces, every $P$-invariant probability on $Y$ has a $P$-invariant lift: the set of lifts is nonempty compact convex, so amenability supplies a fixed point. Consequently stiffness of $X$ implies stiffness of $Y$. The nontrivial flag space $H/P$ is not stiff, because its point mass at $eP$ is $P$-invariant and not $H$-invariant. Hence a compact stiff $H$-space has no such factor when $P$ is amenable. Arbitrary noncompact factor maps need not have the compact lift set, and general parabolic subgroups need not be amenable. These diagnostics do not classify all stiff actions.
\subsection{Short English statement}
The supplied literature reports stiffness for the stated unit-area differential strata. The additional results are diagnostics: extra negative-root symmetries force full invariance, one symmetry outside the $AN$ normalizer suffices for $SL(2,\R)$, and amenability lifts measures through compact factors. The flag-space obstruction uses those compactness and amenability assumptions and is not an unrestricted factor theorem.
\subsection{Sources}
\begin{itemize}
\item Local registry: permanent catalogue entry, exact source code catalogue entry; \texttt{aim-other-notes.json}, \texttt{notes[81]}; canonical tag ``question''.
\item AIM workshop \emph{Emerging applications of measure rigidity} (2004), Question 8 attributed to E. Lindenstrauss, \url{https://aimath.org/WWN/measrigid/measrigid.pdf}.
\item Original reference labels: [64] for stiffness; [48, 199, 141] for the moduli/homogeneous-action background; [143] and [51] for reported partial rigidity results. No complete bibliography entries are supplied.
\item The appended Section 3.2 reports A. Katok's Finsler-sphere geodesic-flow example with two periodic orbits, reference [92]. It is following-section context, not a condition or claim established in the stiffness argument.
\item Supplied research attempt 1, registry updated 14 August 2026: partial result, low novelty confidence, machine-generated and not peer reviewed. It attributes the unit-area stratum theorem to Eskin--Mirzakhani and supplies root/factor diagnostics without complete authoritative citations or orientation-cover details. Verification preserves the extracted record and its section spillover.
\end{itemize}
\end{document}
